% Options for packages loaded elsewhere
\PassOptionsToPackage{unicode}{hyperref}
\PassOptionsToPackage{hyphens}{url}
%
\documentclass[
  12pt,
  a4paper,
  DIV=11,
  numbers=noendperiod]{scrreprt}

\usepackage{amsmath,amssymb}
\usepackage{setspace}
\usepackage{iftex}
\ifPDFTeX
  \usepackage[T1]{fontenc}
  \usepackage[utf8]{inputenc}
  \usepackage{textcomp} % provide euro and other symbols
\else % if luatex or xetex
  \usepackage{unicode-math}
  \defaultfontfeatures{Scale=MatchLowercase}
  \defaultfontfeatures[\rmfamily]{Ligatures=TeX,Scale=1}
\fi
\usepackage{lmodern}
\ifPDFTeX\else  
    % xetex/luatex font selection
\fi
% Use upquote if available, for straight quotes in verbatim environments
\IfFileExists{upquote.sty}{\usepackage{upquote}}{}
\IfFileExists{microtype.sty}{% use microtype if available
  \usepackage[]{microtype}
  \UseMicrotypeSet[protrusion]{basicmath} % disable protrusion for tt fonts
}{}
\usepackage{xcolor}
\setlength{\emergencystretch}{3em} % prevent overfull lines
\setcounter{secnumdepth}{5}
% Make \paragraph and \subparagraph free-standing
\makeatletter
\ifx\paragraph\undefined\else
  \let\oldparagraph\paragraph
  \renewcommand{\paragraph}{
    \@ifstar
      \xxxParagraphStar
      \xxxParagraphNoStar
  }
  \newcommand{\xxxParagraphStar}[1]{\oldparagraph*{#1}\mbox{}}
  \newcommand{\xxxParagraphNoStar}[1]{\oldparagraph{#1}\mbox{}}
\fi
\ifx\subparagraph\undefined\else
  \let\oldsubparagraph\subparagraph
  \renewcommand{\subparagraph}{
    \@ifstar
      \xxxSubParagraphStar
      \xxxSubParagraphNoStar
  }
  \newcommand{\xxxSubParagraphStar}[1]{\oldsubparagraph*{#1}\mbox{}}
  \newcommand{\xxxSubParagraphNoStar}[1]{\oldsubparagraph{#1}\mbox{}}
\fi
\makeatother


\providecommand{\tightlist}{%
  \setlength{\itemsep}{0pt}\setlength{\parskip}{0pt}}\usepackage{longtable,booktabs,array}
\usepackage{calc} % for calculating minipage widths
% Correct order of tables after \paragraph or \subparagraph
\usepackage{etoolbox}
\makeatletter
\patchcmd\longtable{\par}{\if@noskipsec\mbox{}\fi\par}{}{}
\makeatother
% Allow footnotes in longtable head/foot
\IfFileExists{footnotehyper.sty}{\usepackage{footnotehyper}}{\usepackage{footnote}}
\makesavenoteenv{longtable}
\usepackage{graphicx}
\makeatletter
\newsavebox\pandoc@box
\newcommand*\pandocbounded[1]{% scales image to fit in text height/width
  \sbox\pandoc@box{#1}%
  \Gscale@div\@tempa{\textheight}{\dimexpr\ht\pandoc@box+\dp\pandoc@box\relax}%
  \Gscale@div\@tempb{\linewidth}{\wd\pandoc@box}%
  \ifdim\@tempb\p@<\@tempa\p@\let\@tempa\@tempb\fi% select the smaller of both
  \ifdim\@tempa\p@<\p@\scalebox{\@tempa}{\usebox\pandoc@box}%
  \else\usebox{\pandoc@box}%
  \fi%
}
% Set default figure placement to htbp
\def\fps@figure{htbp}
\makeatother

\usepackage{indentfirst}
\usepackage{float}
\floatplacement{figure}{H}
\IfFileExists{plex-otf.sty}{
  %% Full TeXlive
  \usepackage[
    % math,
    RM={Scale=0.94},SS={Scale=0.94},SScon={Scale=0.94},TT={Scale=MatchLowercase,FakeStretch=0.9},
    DefaultFeatures={Ligatures=Common}
  ]{plex-otf}
  % \usepackage[Scale=MatchUppercase]{juliamono}
}{
  %% TinyTeX
  \usepackage{libertine}
}
\KOMAoption{captions}{tableheading}
\makeatletter
\@ifpackageloaded{caption}{}{\usepackage{caption}}
\AtBeginDocument{%
\ifdefined\contentsname
  \renewcommand*\contentsname{Содержание}
\else
  \newcommand\contentsname{Содержание}
\fi
\ifdefined\listfigurename
  \renewcommand*\listfigurename{Список иллюстраций}
\else
  \newcommand\listfigurename{Список иллюстраций}
\fi
\ifdefined\listtablename
  \renewcommand*\listtablename{Список таблиц}
\else
  \newcommand\listtablename{Список таблиц}
\fi
\ifdefined\figurename
  \renewcommand*\figurename{Рисунок}
\else
  \newcommand\figurename{Рисунок}
\fi
\ifdefined\tablename
  \renewcommand*\tablename{Таблица}
\else
  \newcommand\tablename{Таблица}
\fi
}
\@ifpackageloaded{float}{}{\usepackage{float}}
\floatstyle{ruled}
\@ifundefined{c@chapter}{\newfloat{codelisting}{h}{lop}}{\newfloat{codelisting}{h}{lop}[chapter]}
\floatname{codelisting}{Список}
\newcommand*\listoflistings{\listof{codelisting}{Листинги}}
\makeatother
\makeatletter
\makeatother
\makeatletter
\@ifpackageloaded{caption}{}{\usepackage{caption}}
\@ifpackageloaded{subcaption}{}{\usepackage{subcaption}}
\makeatother

\ifLuaTeX
\usepackage[bidi=basic]{babel}
\else
\usepackage[bidi=default]{babel}
\fi
\babelprovide[main,import]{russian}
\babelprovide[import]{english}
% get rid of language-specific shorthands (see #6817):
\let\LanguageShortHands\languageshorthands
\def\languageshorthands#1{}
\usepackage[style=gost-numeric,backend=biber,langhook=extras,autolang=other*]{biblatex}
\addbibresource{bib/cite.bib}
\usepackage{csquotes}
\usepackage{bookmark}

\IfFileExists{xurl.sty}{\usepackage{xurl}}{} % add URL line breaks if available
\urlstyle{same} % disable monospaced font for URLs
\hypersetup{
  pdftitle={Отчёт по лабораторной работе №5},
  pdfauthor={Курилко-Рюмин Евгений Михайлович},
  pdflang={ru-RU},
  hidelinks,
  pdfcreator={LaTeX via pandoc}}


\title{Отчёт по лабораторной работе №5}
\usepackage{etoolbox}
\makeatletter
\providecommand{\subtitle}[1]{% add subtitle to \maketitle
  \apptocmd{\@title}{\par {\large #1 \par}}{}{}
}
\makeatother
\subtitle{Аппарат сетей Петри: задача обедающих философов}
\author{Курилко-Рюмин Евгений Михайлович}
\date{}

\begin{document}
\maketitle

\renewcommand*\contentsname{Содержание}
{
\setcounter{tocdepth}{1}
\tableofcontents
}
\listoffigures
\listoftables

\setstretch{1.5}
\chapter{Цель
работы}\label{ux446ux435ux43bux44c-ux440ux430ux431ux43eux442ux44b}

Целью работы является построение и исследование модели задачи обедающих
философов в аппарате сетей Петри, сравнение классической схемы и
варианта с арбитром, а также подготовка воспроизводимых программных
материалов и отчёта по результатам вычислительных экспериментов.

\chapter{Задание}\label{ux437ux430ux434ux430ux43dux438ux435}

В ходе лабораторной работы требовалось:

\begin{enumerate}
\def\labelenumi{\arabic{enumi}.}
\tightlist
\item
  Подготовить проект \texttt{Julia} в структуре \texttt{DrWatson}.
\item
  Реализовать модель сети Петри для задачи обедающих философов.
\item
  Сформировать literate-версии основных сценариев.
\item
  Выполнить базовый вычислительный эксперимент для двух топологий сети.
\item
  Провести серию запусков для набора параметров.
\item
  Построить графики, итоговый сравнительный отчёт и анимацию.
\item
  Включить результаты выполнения в итоговый документ.
\end{enumerate}

\chapter{Теоретическое
введение}\label{ux442ux435ux43eux440ux435ux442ux438ux447ux435ux441ux43aux43eux435-ux432ux432ux435ux434ux435ux43dux438ux435}

Сеть Петри представляет собой ориентированную двудольную структуру, в
которой позиции описывают состояния системы, переходы моделируют
события, а маркировка задаёт текущее распределение фишек по позициям.
Такой аппарат удобен для описания параллельных процессов, конфликтов за
ресурсы и ситуаций взаимной блокировки \autocite{murata1989petri}.

Задача обедающих философов является классическим примером конкуренции за
разделяемые ресурсы и потому естественно формулируется средствами сетей
Петри \autocite{korolkova2026simulation}. Для каждого философа в модели
используются состояния \texttt{Think\_i}, \texttt{Hungry\_i} и
\texttt{Eat\_i}, а для каждой вилки вводится позиция \texttt{Fork\_i}. В
классической сети каждый философ может одновременно перейти к захвату
первой вилки, из-за чего система способна попасть в тупик. В
модифицированной схеме добавляется позиция \texttt{Arbiter} с
\texttt{N\ -\ 1} фишками; она ограничивает число одновременно
претендующих на ресурсы философов и тем самым разрушает конфигурацию
deadlock.

В реализованном проекте рассматриваются две формы динамики:

\begin{itemize}
\tightlist
\item
  стохастическая симуляция по схеме Гиллеспи;
\item
  детерминированная \texttt{ODE}-аппроксимация как дополнительный
  аналитический слой.
\end{itemize}

\chapter{Выполнение лабораторной
работы}\label{ux432ux44bux43fux43eux43bux43dux435ux43dux438ux435-ux43bux430ux431ux43eux440ux430ux442ux43eux440ux43dux43eux439-ux440ux430ux431ux43eux442ux44b}

\section{Подготовка среды и исходных
материалов}\label{ux43fux43eux434ux433ux43eux442ux43eux432ux43aux430-ux441ux440ux435ux434ux44b-ux438-ux438ux441ux445ux43eux434ux43dux44bux445-ux43cux430ux442ux435ux440ux438ux430ux43bux43eux432}

Работа велась в каталоге \texttt{labs/lab05/project}. В начале была
запущена среда \texttt{Julia}, после чего активировано окружение проекта
и подготовлены зависимости, необходимые для моделирования, визуализации
и генерации производных материалов.

\begin{figure}

\centering{

\includegraphics[width=0.98\linewidth,height=\textheight,keepaspectratio]{image/1.png}

}

\caption{\label{fig-lab05-julia}Запуск интерпретатора Julia в каталоге
проекта}

\end{figure}%

Структура лабораторной работы организована вокруг нескольких ключевых
файлов.

\begin{longtable}[]{@{}
  >{\raggedright\arraybackslash}p{(\linewidth - 2\tabcolsep) * \real{0.5000}}
  >{\raggedright\arraybackslash}p{(\linewidth - 2\tabcolsep) * \real{0.5000}}@{}}
\caption{Основные артефакты лабораторной
работы}\label{tbl-lab05-files}\tabularnewline
\toprule\noalign{}
\begin{minipage}[b]{\linewidth}\raggedright
Файл
\end{minipage} & \begin{minipage}[b]{\linewidth}\raggedright
Назначение
\end{minipage} \\
\midrule\noalign{}
\endfirsthead
\toprule\noalign{}
\begin{minipage}[b]{\linewidth}\raggedright
Файл
\end{minipage} & \begin{minipage}[b]{\linewidth}\raggedright
Назначение
\end{minipage} \\
\midrule\noalign{}
\endhead
\bottomrule\noalign{}
\endlastfoot
\texttt{src/DiningPhilosophers.jl} & описание сети Петри, функции
симуляции и анализа \\
\texttt{scripts/01\_dining\_philosophers.jl} & базовый сценарий
сравнения двух сетей \\
\texttt{scripts/02\_dining\_philosophers\_param.jl} & серия прогонов для
набора параметров \\
\texttt{scripts/dining\_philosophers\_report.jl} & построение итогового
сравнительного графика \\
\texttt{scripts/dining\_philosophers\_animation.jl} & генерация анимации
изменения маркировки \\
\texttt{scripts/generate\_all.jl} & создание \texttt{clean},
\texttt{qmd} и \texttt{ipynb} производных файлов \\
\end{longtable}

После подготовки исходников были сгенерированы производные форматы.

\begin{figure}

\centering{

\includegraphics[width=0.98\linewidth,height=\textheight,keepaspectratio]{image/2.png}

}

\caption{\label{fig-lab05-generate-start}Начало генерации
literate-материалов для базового и параметрического сценариев}

\end{figure}%

\begin{figure}

\centering{

\includegraphics[width=0.98\linewidth,height=\textheight,keepaspectratio]{image/8.png}

}

\caption{\label{fig-lab05-generate-end}Завершение генерации производных
\texttt{jl}, \texttt{qmd} и \texttt{ipynb} файлов}

\end{figure}%

Из этих скриншотов видно, что из literate-скриптов были автоматически
получены исполняемые файлы, \texttt{Quarto}-документы и
\texttt{Jupyter\ notebook}, которые затем использовались как единый
воспроизводимый комплект материалов.

\section{Базовый вычислительный
эксперимент}\label{ux431ux430ux437ux43eux432ux44bux439-ux432ux44bux447ux438ux441ux43bux438ux442ux435ux43bux44cux43dux44bux439-ux44dux43aux441ux43fux435ux440ux438ux43cux435ux43dux442}

Базовый сценарий был реализован в файле
\texttt{scripts/01\_dining\_philosophers.jl}. Для прямого сравнения двух
топологий использовались одинаковые параметры:

\begin{itemize}
\tightlist
\item
  число философов \texttt{N\ =\ 5};
\item
  горизонт моделирования \texttt{tmax\ =\ 80};
\item
  скорость захвата первой вилки \texttt{left\_rate\ =\ 0.5};
\item
  скорость захвата второй вилки \texttt{right\_rate\ =\ 1.8};
\item
  скорость возврата вилок \texttt{put\_rate\ =\ 0.8};
\item
  одинаковое начальное зерно генератора случайных чисел.
\end{itemize}

Запуск основного сценария и полученная сводка показаны на
рис.~\ref{fig-lab05-basic-run} и рис.~\ref{fig-lab05-basic-summary}.

\begin{figure}

\centering{

\includegraphics[width=0.98\linewidth,height=\textheight,keepaspectratio]{image/3.png}

}

\caption{\label{fig-lab05-basic-run}Запуск основного сценария
\texttt{01\_dining\_philosophers.jl}}

\end{figure}%

\begin{figure}

\centering{

\includegraphics[width=0.98\linewidth,height=\textheight,keepaspectratio]{image/5.png}

}

\caption{\label{fig-lab05-basic-summary}Сводные показатели базового
сравнения двух сетей}

\end{figure}%

Численные итоги базового прогона сведены в табл.~\ref{tbl-lab05-basic}.

\begin{longtable}[]{@{}
  >{\raggedright\arraybackslash}p{(\linewidth - 10\tabcolsep) * \real{0.1364}}
  >{\raggedright\arraybackslash}p{(\linewidth - 10\tabcolsep) * \real{0.1364}}
  >{\raggedleft\arraybackslash}p{(\linewidth - 10\tabcolsep) * \real{0.1818}}
  >{\raggedleft\arraybackslash}p{(\linewidth - 10\tabcolsep) * \real{0.1818}}
  >{\raggedleft\arraybackslash}p{(\linewidth - 10\tabcolsep) * \real{0.1818}}
  >{\raggedleft\arraybackslash}p{(\linewidth - 10\tabcolsep) * \real{0.1818}}@{}}
\caption{Сводка базового вычислительного
эксперимента}\label{tbl-lab05-basic}\tabularnewline
\toprule\noalign{}
\begin{minipage}[b]{\linewidth}\raggedright
Сеть
\end{minipage} & \begin{minipage}[b]{\linewidth}\raggedright
Наличие deadlock
\end{minipage} & \begin{minipage}[b]{\linewidth}\raggedleft
Время завершения траектории
\end{minipage} & \begin{minipage}[b]{\linewidth}\raggedleft
Среднее число едящих
\end{minipage} & \begin{minipage}[b]{\linewidth}\raggedleft
Максимум \texttt{Eat}
\end{minipage} & \begin{minipage}[b]{\linewidth}\raggedleft
Число голодных в финале
\end{minipage} \\
\midrule\noalign{}
\endfirsthead
\toprule\noalign{}
\begin{minipage}[b]{\linewidth}\raggedright
Сеть
\end{minipage} & \begin{minipage}[b]{\linewidth}\raggedright
Наличие deadlock
\end{minipage} & \begin{minipage}[b]{\linewidth}\raggedleft
Время завершения траектории
\end{minipage} & \begin{minipage}[b]{\linewidth}\raggedleft
Среднее число едящих
\end{minipage} & \begin{minipage}[b]{\linewidth}\raggedleft
Максимум \texttt{Eat}
\end{minipage} & \begin{minipage}[b]{\linewidth}\raggedleft
Число голодных в финале
\end{minipage} \\
\midrule\noalign{}
\endhead
\bottomrule\noalign{}
\endlastfoot
Классическая & \texttt{true} & 46.33 & 0.675 & 2.0 & 5 \\
С арбитром & \texttt{false} & 79.63 & 0.603 & 2.0 & 3 \\
\end{longtable}

Далее были построены графики эволюции маркировки по основным группам
состояний для классической и модифицированной сети.

\begin{figure}

\centering{

\includegraphics[width=0.92\linewidth,height=\textheight,keepaspectratio]{../project/plots/01_dining_philosophers/classic_simulation.png}

}

\caption{\label{fig-lab05-classic}Эволюция маркировки в классической
сети Петри}

\end{figure}%

\begin{figure}

\centering{

\includegraphics[width=0.92\linewidth,height=\textheight,keepaspectratio]{../project/plots/01_dining_philosophers/arbiter_simulation.png}

}

\caption{\label{fig-lab05-arbiter}Эволюция маркировки в сети Петри с
арбитром}

\end{figure}%

На рис.~\ref{fig-lab05-classic} видно, что в классической постановке
система сначала демонстрирует чередование состояний, однако к концу
прогона все философы оказываются в режиме ожидания, а число свободных
вилок обнуляется. На рис.~\ref{fig-lab05-arbiter} динамика остаётся
живой на всём интервале: позиция \texttt{Arbiter} не допускает
одновременного входа всех философов в опасную конфигурацию.

Для более компактного сравнения использовались суммарные характеристики
по состоянию \texttt{Eat}.

\begin{figure}

\centering{

\includegraphics[width=0.9\linewidth,height=\textheight,keepaspectratio]{../project/plots/01_dining_philosophers/dining_eating_compare.png}

}

\caption{\label{fig-lab05-eat-compare}Сравнение стохастической динамики
числа едящих философов}

\end{figure}%

\begin{figure}

\centering{

\includegraphics[width=0.9\linewidth,height=\textheight,keepaspectratio]{../project/plots/01_dining_philosophers/dining_ode_compare.png}

}

\caption{\label{fig-lab05-ode}Детерминированная аппроксимация суммарного
состояния \texttt{Eat}}

\end{figure}%

Стохастический график на рис.~\ref{fig-lab05-eat-compare} показывает
принципиальное расхождение сценариев: в классической сети кривая уходит
к нулю и перестаёт меняться, тогда как в сети с арбитром сохраняются
колебания. \texttt{ODE}-аппроксимация на рис.~\ref{fig-lab05-ode}
воспроизводит ту же качественную тенденцию и подтверждает более высокую
устойчивость модифицированной модели.

\section{Серия запусков для набора
параметров}\label{ux441ux435ux440ux438ux44f-ux437ux430ux43fux443ux441ux43aux43eux432-ux434ux43bux44f-ux43dux430ux431ux43eux440ux430-ux43fux430ux440ux430ux43cux435ux442ux440ux43eux432}

Вторая часть работы была посвящена серии стохастических прогонов, где
изменялась скорость захвата первой вилки. Этот параметр напрямую влияет
на то, насколько быстро философы входят в конфликт за ресурсы.

Фиксированными оставались:

\begin{itemize}
\tightlist
\item
  \texttt{N\ =\ 5};
\item
  \texttt{right\_rate\ =\ 1.8};
\item
  \texttt{put\_rate\ =\ 0.8};
\item
  \texttt{tmax\ =\ 80};
\item
  число прогонов для каждого случая \texttt{40}.
\end{itemize}

Перебирался набор значений

\[
left\_rate \in \{0.3,\;0.5,\;0.7,\;0.9,\;1.1,\;1.3\}.
\]

Командный запуск сценария \texttt{02\_dining\_philosophers\_param.jl} и
часть его результатов представлены на рис.~\ref{fig-lab05-param-run} и
рис.~\ref{fig-lab05-param-summary}.

\begin{figure}

\centering{

\includegraphics[width=0.98\linewidth,height=\textheight,keepaspectratio]{image/4.png}

}

\caption{\label{fig-lab05-param-run}Запуск параметрического сценария
\texttt{02\_dining\_philosophers\_param.jl}}

\end{figure}%

\begin{figure}

\centering{

\includegraphics[width=0.98\linewidth,height=\textheight,keepaspectratio]{image/6.png}

}

\caption{\label{fig-lab05-param-summary}Сводная таблица параметрического
исследования}

\end{figure}%

Сводные численные показатели приведены в табл.~\ref{tbl-lab05-param}.

\begin{longtable}[]{@{}
  >{\raggedleft\arraybackslash}p{(\linewidth - 10\tabcolsep) * \real{0.1667}}
  >{\raggedleft\arraybackslash}p{(\linewidth - 10\tabcolsep) * \real{0.1667}}
  >{\raggedleft\arraybackslash}p{(\linewidth - 10\tabcolsep) * \real{0.1667}}
  >{\raggedleft\arraybackslash}p{(\linewidth - 10\tabcolsep) * \real{0.1667}}
  >{\raggedleft\arraybackslash}p{(\linewidth - 10\tabcolsep) * \real{0.1667}}
  >{\raggedleft\arraybackslash}p{(\linewidth - 10\tabcolsep) * \real{0.1667}}@{}}
\caption{Результаты серии запусков для набора
параметров}\label{tbl-lab05-param}\tabularnewline
\toprule\noalign{}
\begin{minipage}[b]{\linewidth}\raggedleft
\texttt{left\_rate}
\end{minipage} & \begin{minipage}[b]{\linewidth}\raggedleft
Вероятность deadlock, классическая
\end{minipage} & \begin{minipage}[b]{\linewidth}\raggedleft
Среднее время deadlock, классическая
\end{minipage} & \begin{minipage}[b]{\linewidth}\raggedleft
Среднее число едящих, классическая
\end{minipage} & \begin{minipage}[b]{\linewidth}\raggedleft
Вероятность deadlock, арбитр
\end{minipage} & \begin{minipage}[b]{\linewidth}\raggedleft
Среднее число едящих, арбитр
\end{minipage} \\
\midrule\noalign{}
\endfirsthead
\toprule\noalign{}
\begin{minipage}[b]{\linewidth}\raggedleft
\texttt{left\_rate}
\end{minipage} & \begin{minipage}[b]{\linewidth}\raggedleft
Вероятность deadlock, классическая
\end{minipage} & \begin{minipage}[b]{\linewidth}\raggedleft
Среднее время deadlock, классическая
\end{minipage} & \begin{minipage}[b]{\linewidth}\raggedleft
Среднее число едящих, классическая
\end{minipage} & \begin{minipage}[b]{\linewidth}\raggedleft
Вероятность deadlock, арбитр
\end{minipage} & \begin{minipage}[b]{\linewidth}\raggedleft
Среднее число едящих, арбитр
\end{minipage} \\
\midrule\noalign{}
\endhead
\bottomrule\noalign{}
\endlastfoot
0.3 & 0.325 & 41.88 & 0.641 & 0.0 & 0.673 \\
0.5 & 0.800 & 33.20 & 0.565 & 0.0 & 0.618 \\
0.7 & 0.975 & 18.77 & 0.501 & 0.0 & 0.543 \\
0.9 & 1.000 & 15.05 & 0.431 & 0.0 & 0.507 \\
1.1 & 1.000 & 9.94 & 0.419 & 0.0 & 0.479 \\
1.3 & 1.000 & 7.44 & 0.378 & 0.0 & 0.459 \\
\end{longtable}

\begin{figure}

\centering{

\includegraphics[width=0.92\linewidth,height=\textheight,keepaspectratio]{../project/plots/02_dining_philosophers_param/dining_param_panel.png}

}

\caption{\label{fig-lab05-param-panel}Сводный графический отчёт по
набору параметров}

\end{figure}%

По табл.~\ref{tbl-lab05-param} и рис.~\ref{fig-lab05-param-panel} можно
сделать три основных вывода:

\begin{enumerate}
\def\labelenumi{\arabic{enumi}.}
\tightlist
\item
  При росте \texttt{left\_rate} классическая сеть всё быстрее переходит
  в тупиковое состояние.
\item
  Одновременно уменьшается средняя активность системы, то есть число
  философов в состоянии \texttt{Eat}.
\item
  Сеть с арбитром сохраняет нулевую вероятность deadlock на всём
  исследованном диапазоне значений.
\end{enumerate}

\section{Постобработка результатов и итоговые
артефакты}\label{ux43fux43eux441ux442ux43eux431ux440ux430ux431ux43eux442ux43aux430-ux440ux435ux437ux443ux43bux44cux442ux430ux442ux43eux432-ux438-ux438ux442ux43eux433ux43eux432ux44bux435-ux430ux440ux442ux435ux444ux430ux43aux442ux44b}

После завершения вычислительных экспериментов выполнялась отдельная
постобработка сохранённых данных. Для этого использовался скрипт
\texttt{scripts/dining\_philosophers\_report.jl}, который строил
итоговый сравнительный график по состояниям \texttt{Eat\_i} на основе
ранее созданных \texttt{CSV}-файлов.

\begin{figure}

\centering{

\includegraphics[width=0.98\linewidth,height=\textheight,keepaspectratio]{image/7.png}

}

\caption{\label{fig-lab05-final-script}Запуск скрипта построения
итогового графика}

\end{figure}%

\begin{figure}

\centering{

\includegraphics[width=0.9\linewidth,height=\textheight,keepaspectratio]{../project/plots/dining_philosophers_report/final_report.png}

}

\caption{\label{fig-lab05-final}Итоговый сравнительный график по
состояниям \texttt{Eat\_i}}

\end{figure}%

Именно этот график наиболее наглядно разделяет два сценария. Для
классической сети кривые довольно быстро замирают на нуле, что отражает
потерю живости. Для сети с арбитром сохраняются повторяющиеся пики, то
есть возможность приёма пищи остаётся доступной до конца моделирования.

Кроме того, был подготовлен анимационный файл
\texttt{plots/dining\_philosophers\_animation/philosophers\_simulation.gif}.
Он не используется как основной источник количественного анализа, но
удобен для визуального просмотра того момента, когда классическая сеть
прекращает изменять маркировку.

\section{Обсуждение полученных
результатов}\label{ux43eux431ux441ux443ux436ux434ux435ux43dux438ux435-ux43fux43eux43bux443ux447ux435ux43dux43dux44bux445-ux440ux435ux437ux443ux43bux44cux442ux430ux442ux43eux432}

По сравнению с краткой схемой выполнения лабораторной работы здесь были
разнесены три уровня анализа:

\begin{itemize}
\tightlist
\item
  фиксированный базовый прогон, позволяющий увидеть одиночную
  траекторию;
\item
  серия запусков для набора параметров, показывающая устойчивость
  наблюдаемого эффекта;
\item
  отдельная постобработка, в которой результаты сводятся к удобным
  итоговым графикам и артефактам документации.
\end{itemize}

Такое представление отличается от эталонного образца по композиции,
однако сохраняет его академическую логику: сначала показывается
подготовка среды, затем ход эксперимента, после чего отдельно
обсуждаются результаты и выводы. Добавленные терминальные скриншоты
фиксируют реальное выполнение сценариев, генерацию literate-материалов и
сборку итоговых графиков.

\chapter{Выводы}\label{ux432ux44bux432ux43eux434ux44b}

В ходе лабораторной работы задача обедающих философов была успешно
смоделирована средствами сетей Петри, а результаты оформлены в виде
воспроизводимого вычислительного проекта.

\begin{enumerate}
\def\labelenumi{\arabic{enumi}.}
\tightlist
\item
  Реализованы классическая сеть и её модификация с арбитром.
\item
  Подготовлены literate-скрипты, исполняемые файлы,
  \texttt{Jupyter\ notebook} и \texttt{Quarto}-документы.
\item
  Базовый сценарий показал возникновение deadlock в классической сети
  при сохранении работоспособности модели с арбитром.
\item
  Серия запусков подтвердила, что увеличение скорости захвата первой
  вилки резко ухудшает устойчивость классической сети, тогда как наличие
  арбитра удерживает вероятность тупика на нуле.
\item
  Итоговые графики, таблицы, скриншоты и анимация позволяют считать
  лабораторную работу полностью воспроизводимой и документированной.
\end{enumerate}

\chapter*{Список
литературы}\label{ux441ux43fux438ux441ux43eux43a-ux43bux438ux442ux435ux440ux430ux442ux443ux440ux44b}
\addcontentsline{toc}{chapter}{Список литературы}

\printbibliography[heading=none]





\end{document}
